\chapter{جمع‌بندی و پیشنهادها}
%=====================================================================

فناوری توزیع کلید کوانتومی و رمزنگاری پساکوانتومی دیگر در مرحلهٔ آزمایشگاهی نیستند؛ همان‌طور که در فصل‌های \lr{6} و \lr{7} نشان داده شد، بانک‌های پیشرو جهان و کشورهایی نظیر چین، سوئد و سنگاپور، این فناوری‌ها را در محیط‌های عملیاتی واقعی به‌کار گرفته‌اند. با نهایی‌شدن استانداردهای \lr{NIST} در سال $2024$، مسیر مهاجرت به رمزنگاری پساکوانتومی از نظر فنی کاملاً روشن و قابل‌اجراست، و فناوری \lr{QKD} نیز -- به‌ویژه برای کاربردهای کلان‌شهری با فاصلهٔ کوتاه نظیر اتصال مراکز دادهٔ بانکی -- به بلوغ تجاری کافی رسیده است.

این پروپوزال پیشنهاد می‌کند که جمهوری اسلامی ایران، به‌عنوان گامی پیشدستانه در برابر تهدید «برداشت اکنون، رمزگشایی بعداً»، یک طرح پایلوت کلان‌شهری در تهران را با تمرکز بر اتصال بانک مرکزی و زیرساخت‌های تسویهٔ بانکی، هم‌زمان با آغاز فوری مهاجرت نرم‌افزاری به \lr{ML-KEM} و \lr{ML-DSA} در سطح کل شبکهٔ بانکی کشور، آغاز کند. اجرای موفق این طرح پایلوت -- با بودجهٔ برآوردی $3.8$ تا $6.9$ میلیون دلار و افق زمانی $12$ ماهه -- می‌تواند بستری برای توسعهٔ دانش فنی بومی، کاهش وابستگی راهبردی به فناوری خارجی، و ارتقای جایگاه ایران در حوزهٔ نوظهور امنیت کوانتومی منطقه فراهم آورد.

\textbf{پیشنهاد نهایی جهت اقدام:} تشکیل یک کارگروه مشترک میان بانک مرکزی، معاونت علمی، فناوری و اقتصاد دانش‌بنیان ریاست‌جمهوری، مرکز افتا، و شرکت مخابرات ایران برای آغاز فاز \lr{1} (امکان‌سنجی) ظرف مدت $60$ روز از تصویب این طرح.

